\documentclass[10pt,a4paper]{amsart}
\usepackage[margin=19mm]{geometry}
\usepackage{amsmath,amssymb,mathtools}
\usepackage[scr=rsfs,cal=euler]{mathalfa}
\usepackage{xfrac}
\usepackage[unicode,hidelinks,bookmarks=false]{hyperref}
\newcommand{\ie}{i.e.\ }
\newcommand{\cf}{cf.\ }
\newcommand{\resp}{respectively\ }
\newcommand{\Subsection}{\subsection}
\providecommand{\nobreakdash}{\nobreak\hbox{-}}
\setlength{\emergencystretch}{2em}
\multlinegap0pt
\usepackage{mdframed}
\usepackage{mdframed}
\usepackage{mdframed}
\usepackage{mdframed}
\usepackage{mdframed}
\usepackage{csquotes}
\usepackage{amssymb}
\usepackage{dsfont}
\usepackage{xcolor}
\usepackage{braket}
\usepackage[shortlabels]{enumitem}
\usepackage{mathtools}
\usepackage{cancel}
\newtheorem{cor}{}
\newcommand{\Id}{\mathds{1}}
\newcommand{\Ord}[1]{\calO\left(#1\right)}
\newcommand{\RR}{\mathbb{R}}
\newcommand{\abs}[1]{\left\lvert#1\right\rvert}
\newcommand{\br}[1]{\left(#1\right)}
\newcommand{\calO}{\mathcal{O}}
\newcommand{\calR}{\mathcal{R}}
\newcommand{\dif}{\operatorname{d}\!} % shorten the space afterwards
\newcommand{\ee}{\operatorname{e}}
\newcommand{\eq}[1]{\begin{align*}#1\end{align*}}
\newcommand{\eql}[1]{\begin{align}#1\end{align}}
\newcommand{\tr}{\operatorname{tr}}
\title{Instantons in a Double-Well are Poisson Distributed}
\author{Jacob Shapiro}
\date{Source: arXiv:2608.23342v1 (CC BY 4.0)}
\begin{document}
\maketitle

\noindent\emph{Source and licence.} Instantons in a Double-Well are Poisson Distributed. Version arXiv:2608.23342v1 (CC BY 4.0), distributed by arXiv under the Creative Commons Attribution 4.0 International licence, http://creativecommons.org/licenses/by/4.0/, as recorded in the arXiv licence metadata for this exact version and shown on the abstract page https://arxiv.org/abs/2608.23342v1. Full licence notice: \texttt{licenses/arxiv-2608.23342-CC-BY-4.0.txt}.\par\smallskip

\noindent\textit{Editorial scope.} Counted unit: the article's cor rho (source0.tex lines 810--815). The article's complete original proof of it is delivered with it (source0.tex lines 817--1056, one \texttt{proof} environment of 240 lines). No mathematics is written, repaired or re-derived: the statement and the proof are the article's own lines, in the article's own notation, and no retained line is substituted or re-flowed.  No further source-local context is needed: every cross-reference of every retained line resolves inside the two delivered spans.  Nothing was omitted from any retained span, so the package carries no omission record.  No retained line contains a citation, so the wrapper carries no bibliography and no bibliography entry.  No retained line contains a figure environment, an image-inclusion command or drawing code, so no image is delivered and the package declares no asset dependency.  Editorial formatting only, in addition: an AMS \texttt{amsart} preamble that restates the article's own declarations and macros used by the retained lines, namely the environments \texttt{cor} on separate counters, together with the standard packages those lines need.  The retained environments print in this document as Corollary 1 in order of appearance, whereas the article prints the same items with the numbering of its own full document.  The complete native source of the article as distributed in the arXiv e-print for this exact version is bundled unchanged as \texttt{sources/208407/source0.tex}.
\begin{cor}\label{cor:eigenvalue splitting controlled by rho}
    For the $\rho_\lambda$ defined above, we have
    \eql{
    \lim_{\lambda\to\infty}\frac{E_1(\lambda)-E_0(\lambda)}{2\abs{\rho_\lambda}} = 1\,.
    }
\end{cor}

\begin{proof}
The parity of the endpoint implies the exact decompositions
\eq{
Z_\lambda(\beta)
&=
\sum_{m=0}^{\infty}Z_{\lambda,2m}(\beta),\\
Z_\lambda^{\curvearrowright}(\beta)
&=
\sum_{m=0}^{\infty}Z_{\lambda,2m+1}(\beta).
}
Hence, upon setting
\eq{
\beta_{\lambda,N}:=\frac{N}{2\abs{\rho_\lambda}},
}
we obtain
\eq{
A_\lambda(N)
&:=
\frac{
Z_\lambda(\beta_{\lambda,N})
-
Z_\lambda^{\curvearrowright}(\beta_{\lambda,N})
}{
Z_\lambda(\beta_{\lambda,N})
+
Z_\lambda^{\curvearrowright}(\beta_{\lambda,N})
}\\
&=
\sum_{k=0}^{\infty}
(-1)^k
\mathbb{P}_{\lambda,\beta_{\lambda,N}}^{\mathrm{instanton}}
\left[\Set{k}\right].
}
Applying the theorem, we have
\eq{
\lim_{\lambda\to\infty}
\mathbb{P}_{\lambda,\beta_{\lambda,N}}^{\mathrm{instanton}}
\left[\Set{k}\right]
=
\ee^{-N/2}\frac{(N/2)^k}{k!}.
}
Since the limiting Poisson masses sum to one, this pointwise
convergence implies convergence in total variation (discrete Scheffé lemma). It follows that
\eq{
\lim_{\lambda\to\infty}A_\lambda(N)
&=
\sum_{k=0}^{\infty}
(-1)^k\ee^{-N/2}\frac{(N/2)^k}{k!}\\
&=
\ee^{-N/2}
\sum_{k=0}^{\infty}\frac{(-N/2)^k}{k!}
=
\ee^{-N}.
}

Let
\eq{
\Pi_\pm:=\frac{\Id\pm R}{2}
}
be the projections onto the even and odd subspaces, respectively.
Since $H_\lambda$ and $\chi_\lambda$ commute with $R$, we have
\eq{
Z_\lambda(\beta)\pm Z_\lambda^{\curvearrowright}(\beta)
=
2\tr\br{
\chi_\lambda
\exp\br{-\beta H_\lambda}
\Pi_\pm
\chi_\lambda
}.
}
The range of $P_{\lambda,0}$ lies in the even subspace, whereas the
range of $P_{\lambda,1}$ lies in the odd subspace. Set
\eq{
a_{\lambda,j}
:=
\tr\br{
\chi_\lambda P_{\lambda,j}\chi_\lambda
},
\qquad j\in\Set{0,1}\,,\,\qquad \calR_{\lambda,\pm}(\beta) := \tr\br{\chi_\lambda(X)\exp\br{-\beta H_\lambda}\Pi_\pm \chi_{[E_2(\lambda),\infty)}(H_\lambda) \chi_\lambda(X)}\,.
} %Hence
%\eq{
%\tr\br{
%\chi_\lambda
%\exp\br{-\beta H_\lambda}
%\Pi_\pm
%\chi_\lambda
%} = a_{\lambda,{0,1}}\exp\br{-\beta E_{0,1}(\lambda)}+\calR_{\lambda,\pm}(\beta)
%} so that
Hence
\eq{
A_\lambda(N) = \frac{a_{\lambda,1}
\exp\br{
-N\frac{E_1(\lambda)}
        {2\abs{\rho_\lambda}}
}+\calR_{\lambda,-}(\beta_{\lambda,N})}{a_{\lambda,0}\exp\br{
-N\frac{E_0(\lambda)}
        {2\abs{\rho_\lambda}}
}+\calR_{\lambda,+}(\beta_{\lambda,N})}
}

Now choose e.g.
$\tau_\lambda:=\frac{1}{\lambda^2}$.
Using the spectral theorem, we obtain
\eq{
\calR_{\lambda,\pm}(\beta)
\leq
\exp\br{
-\br{\beta-\tau_\lambda}E_2(\lambda)
}
\tr\br{
\chi_\lambda
\exp\br{-\tau_\lambda H_\lambda}
\chi_\lambda
}.
}
Dividing by the corresponding leading term gives, for
$j\in\Set{0,1}$,
\eq{
\frac{
\calR_{\lambda,\pm}(\beta)
}{
a_{\lambda,j}\exp\br{-\beta E_j(\lambda)}
}
\leq
\frac{
\exp\br{\tau_\lambda E_2(\lambda)}
\tr\br{
\chi_\lambda
\exp\br{-\tau_\lambda H_\lambda}
\chi_\lambda
}
}{
a_{\lambda,j}
}
\exp\br{
-\beta\br{E_2(\lambda)-E_j(\lambda)}
}.
}
At $\beta=\beta_{\lambda,N}$, the relevant suppression factor satisfies
\eq{
\exp\br{
-\beta_{\lambda,N}
\br{E_2(\lambda)-E_1(\lambda)}
}
\leq
\exp\br{
-\frac{N \lambda c_{\rm gap}}{2\abs{\rho_\lambda}}
}.
}
Since $\log(\lambda)\abs{\rho_\lambda}\to0$, this factor tends to zero. To conclude that
the entire right-hand side tends to zero, however, we must also
control the prefactor. By the Feynman--Kac formula and the non-negativity of $V$,
\eq{
\tr\br{
\chi_\lambda
\exp\br{-H_\lambda/\lambda^2}
\chi_\lambda
}
\leq
\br{\frac{\lambda^2}{4\pi}}^{n/2}
\int_{x\in\RR^n}\chi_\lambda(x)^2\dif{x}.
}
Because $\chi_\lambda$ is supported in two balls of radius
$\lambda^{-1/4}$,
\eq{
\int_{x\in\RR^n}\chi_\lambda(x)^2\dif{x}
=
\Ord{\lambda^{-n/4}},
}
and hence
\eq{
\tr\br{
\chi_\lambda
\exp\br{-H_\lambda/\lambda^2}
\chi_\lambda
}
=
\Ord{\lambda^{3n/4}}.
}
We have $E_2(\lambda)\leq\lambda^2 c_{\infty}$ and
$a_{\lambda,j}\gtrsim1-\calO(\lambda^{-1/2})$ by basic energy estimates. The prefactor is therefore at most
polynomial in $\lambda$. Thus, for some $C,M>0$,
\eq{
\frac{
\calR_{\lambda,\pm}(\beta_{\lambda,N})
}{
a_{\lambda,j}
\exp\br{-\beta_{\lambda,N}E_j(\lambda)}
}
\leq
C\lambda^M
\exp\br{
-\frac{N \lambda c_{\rm gap}}{2\abs{\rho_\lambda}}
}
\longrightarrow0\,.
}

We conclude that for every fixed
$N>0$,
\eq{
A_\lambda(N)
=
\frac{a_{\lambda,1}}{a_{\lambda,0}}
\exp\br{
-N\frac{E_1(\lambda)-E_0(\lambda)}
        {2\abs{\rho_\lambda}}
}
\br{1+o(1)}.
}

Applying this asymptotic with $N=1$ and $N=2$ and taking
their ratio eliminates the localization factors:
\eq{
\frac{A_\lambda(2)}{A_\lambda(1)}
=
\exp\br{
-\frac{E_1(\lambda)-E_0(\lambda)}
       {2\abs{\rho_\lambda}}
}
\br{1+o(1)}.
}
On the other hand, the Poisson limit gives
\eq{
\lim_{\lambda\to\infty}
\frac{A_\lambda(2)}{A_\lambda(1)}
=
\frac{\ee^{-2}}{\ee^{-1}}
=
\ee^{-1}.
}
Taking logarithms therefore yields
\eq{
\lim_{\lambda\to\infty}
\frac{E_1(\lambda)-E_0(\lambda)}
     {2\abs{\rho_\lambda}}
=
1\,.
}
\end{proof}

\end{document}
